% =========================================================
% PART 2 — BIAS REGIME AND DETECTION THEOREM (RECONCILED FINAL)
% =========================================================

\section{Bias Regime and Detection Theorem}\label{sec:bias}

\subsection{Setup and Decomposition}

Fix $\sigma\in[1/2,1)$ and let $\rho=\beta+i\gamma$ be a zero of $\zeta(s)$ with
$\eta=\beta-\sigma>0$.  We examine its contribution to the smoothed tail
$H_\sigma(t)$ and to the spline–penalized functional $F_\lambda$
from~\eqref{eq:SPTB}:
\begin{equation*}
H_\sigma(t)
  = \sum_{\rho'} \frac{e^{(\beta'-\sigma)t}}{|\rho'|^{\alpha}}\cos(\gamma' t)
  = h_\rho(t)+R(t),
\qquad
h_\rho(t)=|\rho|^{-\alpha}e^{\eta t}\cos(\gamma t),
\end{equation*}
where $R$ collects the contributions of $\rho'\ne\rho$.

\begin{definition}[Single-zero residual]\label{def:single-zero}
For a block $I_j=[t_j,t_{j+1}]$, the \emph{single-zero residual}
associated to $\rho$ is $h_\rho-S_j$, where $S_j$ is the $L^2(I_j)$-best affine fit to $H_\sigma$.
\end{definition}

% ---------------------------------------------------------
\subsection{Step 1: Derivative Lower Bound}\label{step:derivative-lb}

The derivative term in $F_\lambda$ captures slope mismatch that affine fitting cannot remove.
By Lemma~\ref{lem:derivative-penalty} applied to $h_\rho(t)=|\rho|^{-\alpha}e^{\eta t}\cos(\gamma t)$
on $I_j$, combined with the blockwise derivative--variance constant $c_{\mathrm{der}}=4\pi^2$
(Appendix~\ref{app:B}), we obtain
\begin{equation}
\lambda \int_{I_j} |h_\rho'(t)|^2 dt
  \ge \frac{\lambda\,c_{\mathrm{der}}\,\eta^2}{2|\rho|^{2\alpha}}
      \min\!\Bigl\{\Delta,\frac{1}{|\gamma|}\Bigr\}
      e^{2\eta t_j},
\label{eq:derivative-lb}
\end{equation}
uniformly in $j$.  Summing over blocks introduces a geometric factor in $e^{2\eta t_j}$.

% ---------------------------------------------------------
\subsection{Step 2: Residual Slope Variance}\label{step:residual-variance}

Let $R=H_\sigma-h_\rho$, the partial sum over all zeros $\rho'\ne\rho$.
We bound $\sum_j\int_{I_j}|R'|^2$ in the setting of
Theorem~\ref{thm:bias}, where $\rho$ is the unique off-line zero so that
all remaining zeros $\rho'\ne\rho$ in $R$ satisfy $\beta'\le\sigma$.

Since every $\rho'\ne\rho$ in $R$ satisfies $\beta'\le\sigma$, the residual is a
variance-regime field, and Theorem~\ref{thm:variance} (proved in Part~1 via the
global second moment of Appendix~\ref{app:A}, \emph{not} the large sieve for the
clustered ordinates) applies to it directly:
\begin{equation}
\sum_j\!\int_{I_j}\!|R'(t)|^2dt
   \;\le\; C_0'\,T\log T\log\log T,
\label{eq:residual-variance}
\end{equation}
with $C_0'$ as in~\eqref{eq:residC0prime}.  This bound is essentially
unconditional: by Remark~\ref{rmk:obstruction} it needs only the trivial
low-height ordinate separation, far weaker than Hypothesis~(S).  The cross-term
between $h_\rho$ and $R$ is handled at the global level by the triangle-inequality
assembly of Step~\ref{step:assembly}, so no per-block cancellation is invoked.

% ---------------------------------------------------------
\subsection{Step 3: Affine Projection Residual Energy}\label{step:affine-energy}

Affine projection removes at most the mean of $h_\rho$ on a block; its oscillatory curvature remains.
For $I_j=[t_j,t_{j+1}]$,
\begin{equation}
\int_{I_j}\!|h_\rho-S_j|^2dt
   \gg e^{2\eta t_j}\!\min\!\Bigl\{\Delta,\tfrac{1}{|\gamma|}\Bigr\}.
\label{eq:affine-energy}
\end{equation}
This follows directly from the explicit Gram-matrix computation for
$e^{\eta t}\cos(\gamma t)$ on $[0,\Delta]$ (Lemma~\ref{lem:affine-lb}).

\begin{remark}[Spline smoothing]\label{rmk:spline}
Replacing affine fits by natural $C^2$ cubic splines only rescales constants:
for some absolute $c\in(0,1)$,
\[
\int_I |(h_\rho-S_I^{(3)})'|^2 dt
   \ge c\,\int_I |(h_\rho-S_I)'|^2 dt,
\]
uniformly under $\kappa_1/\log T\le\Delta\le\kappa_2/\log T$
and $\lambda\asymp(\log T)^{-2}$.
\end{remark}

% ---------------------------------------------------------
\subsection{Step 4: Geometric Summation}\label{step:geom-sum}

Summing over blocks gives the geometric series
\begin{equation}
\sum_j e^{2\eta t_j}
   = \frac{e^{2\eta T}-1}{e^{2\eta\Delta}-1}
   \asymp \frac{e^{2\eta T}}{2\eta\Delta},
\label{eq:geom-sum}
\end{equation}
the source of exponential scaling in $T$.

% ---------------------------------------------------------
\subsection{Step 5: Assembly}\label{step:assembly}

Write $H_\sigma=h_\rho+R$ with $R=H_\sigma-h_\rho$, and let $S_j$ be the blockwise
affine projection.  Rather than a per-block subtraction, we argue globally: by the
triangle inequality in $L^2\bigl(\bigcup_j I_j\bigr)$,
\[
\Bigl(\sum_j\!\int_{I_j}\!|\partial_t(H_\sigma-S_j)|^2\Bigr)^{1/2}
\;\ge\;
\Bigl(\sum_j\!\int_{I_j}\!|\partial_t(h_\rho-S_j)|^2\Bigr)^{1/2}
-\Bigl(\sum_j\!\int_{I_j}\!|R'|^2\Bigr)^{1/2},
\]
so the residual cross-term is absorbed by the inequality itself (no
block-by-block cancellation is invoked).  By Steps~\ref{step:derivative-lb}
and~\ref{step:geom-sum} the first term is $\ge E(T)^{1/2}$ with
\[
E(T)=\frac{c_{\mathrm{der}}\,\eta^2}{2|\rho|^{2\alpha}}
        \min\!\left\{\!\Delta,\tfrac{1}{|\gamma|}\!\right\}\frac{e^{2\eta T}}{2\eta\Delta},
\]
and by Step~\ref{step:residual-variance} the second is $\le P(T)^{1/2}$ with
$P(T)=C_0'\,T\log T\log\log T$.  Since $E$ grows exponentially and $P$ only
polynomially, $E(T)\ge 4P(T)$ for all sufficiently large $T$, whence
$E(T)^{1/2}-P(T)^{1/2}\ge \tfrac12 E(T)^{1/2}$ and
\begin{equation}
\lambda\!\sum_j\!\int_{I_j}\!|\partial_t(H_\sigma-S_j)|^2dt
  \;\ge\; \tfrac14\,\lambda\,E(T).
\label{eq:assembly}
\end{equation}
For any fixed $\eta>0$ this is exponential in $T$ and dominates the polynomial
variance regime.

% ---------------------------------------------------------
\subsection{Step 6: Uniform Constants and Threshold Regime}\label{step:uniform}

Tracking constants gives
\begin{equation}
F_\lambda(H_\sigma;T,\Delta)
  \ge \frac{c_{\mathrm{der}}}{4C_0'}\,
        \frac{\lambda\eta^2}{|\rho|^{2\alpha}}\,
        \frac{e^{2\eta T}}{\eta\Delta},
\label{eq:uniform-lb}
\end{equation}
uniformly for $\kappa_1/\log T \le \Delta \le \kappa_2/\log T$ and
$c_1(\log T)^{-2}\le \lambda \le c_2(\log T)^{-2}$,
with $c_{\mathrm{der}}=4\pi^2$ and $C_0'$ the residual variance constant
of~\eqref{eq:residual-variance}.

\begin{lemma}[Threshold regime]\label{lem:threshold}
Fix the canonical scaling and assume $\eta=\beta-\sigma\ge c/\log T$ for constant $c>0$.
Then for all sufficiently large $T$,
\[
F_\lambda(H_\sigma;T,\Delta)
   \ge A(\alpha,\sigma,c)\,\frac{e^{2\eta T}}{\eta\Delta},
\]
for an explicit $A(\alpha,\sigma,c)>0$.
Hence $F_\lambda\ge\exp(c' T/\log T)$ for some $c'=c'(\alpha,\sigma,c)$:
the exponential regime overtakes any polynomial bound.
\end{lemma}

\begin{proof}
By~\eqref{eq:assembly}, the derivative contribution satisfies
\[
\lambda\!\sum_j\!\int_{I_j}\!|\partial_t(H_\sigma-S_j)|^2dt
\;\ge\;
\underbrace{\frac{\lambda\,c_{\mathrm{der}}\,\eta^2}{2|\rho|^{2\alpha}}
\min\!\bigl\{\Delta,\tfrac{1}{|\gamma|}\bigr\}
\frac{e^{2\eta T}}{2\eta\Delta}}_{\displaystyle =:E(T)}
\;-\;
\underbrace{\lambda\,C_0'\,T\log T\log\log T}_{\displaystyle =:P(T)}.
\]
In the canonical regime, $\eta\ge c/\log T$, $\Delta\asymp\kappa/\log T$,
$\lambda\asymp c_\lambda(\log T)^{-2}$, so
$\eta\Delta\asymp c\kappa/\log^2 T$ and
\[
E(T)\;\ge\;
\frac{c_{\mathrm{der}}\,c_\lambda\,c^2\,\kappa}{4|\rho|^{2\alpha}}
\cdot(\log T)^{-4}\cdot
\min\!\bigl\{\tfrac{\kappa}{\log T},\tfrac{1}{|\gamma|}\bigr\}\cdot
\frac{e^{2cT/\log T}}{c\kappa/\log^2 T}.
\]
The exponential factor $e^{2cT/\log T}$ grows faster than any polynomial in~$T$.
Specifically, for $T\ge T_0(\alpha,\sigma,c)$ we have
$E(T)\ge 2P(T)$, since $e^{2cT/\log T}\gg T^k$ for every fixed~$k$.
Hence $F_\lambda\ge E(T)-P(T)\ge E(T)/2$, giving
\[
F_\lambda\;\ge\; A(\alpha,\sigma,c)\,\frac{e^{2\eta T}}{\eta\Delta}
\]
with $A=c_{\mathrm{der}}\,c_\lambda\,c^2/(8|\rho|^{2\alpha})
\cdot\min\{1,(|\gamma|\Delta)^{-1}\}>0$.
Since $\eta\ge c/\log T$ and $\Delta\asymp 1/\log T$, we obtain
$F_\lambda\ge\exp(c'T/\log T)$ for
$c'=2c-\varepsilon$ for any $\varepsilon>0$ and $T$ sufficiently large.
\end{proof}

% ---------------------------------------------------------
\subsection{Main Theorem (Bias/Detection)}\label{subsec:bias-theorem}

\begin{theorem}[Bias regime and detection]\label{thm:bias}
Let $\alpha\ge1$, $\sigma\in[1/2,1)$, and assume the canonical regime
$\kappa_1/\log T \le \Delta \le \kappa_2/\log T$ and
$c_1(\log T)^{-2}\le \lambda \le c_2(\log T)^{-2}$.
If a zero $\rho=\beta+i\gamma$ satisfies $\eta:=\beta-\sigma>0$, then
\[
F_\lambda(H_\sigma;T,\Delta)
   \gg \lambda\,
       \frac{\eta^{2}}{|\rho|^{2\alpha}}\,
       \min\!\Bigl\{\Delta,\tfrac{1}{|\gamma|}\Bigr\}\,
       \frac{e^{2\eta T}}{\eta\Delta},
\]
with an implied constant depending only on
$(\alpha,\sigma,\kappa_1,\kappa_2,c_1,c_2)$.
In particular, any $\beta>\sigma$ forces exponential growth with slope $2(\beta-\sigma)$.
\end{theorem}

\begin{proof}
Immediate from Lemma~\ref{lem:threshold} and the uniform lower
bound~\eqref{eq:uniform-lb}: the single-zero derivative-penalty contribution
(Steps~\ref{step:derivative-lb}--\ref{step:assembly}) dominates the residual
slope variance~\eqref{eq:residual-variance} once $e^{2\eta T}$ exceeds the
polynomial bound $C_0\,T\log T\log\log T$, which holds for all sufficiently
large~$T$ in the canonical regime.
\end{proof}

\begin{remark}[Multiple off-line zeros]\label{rmk:multiple-zeros}
When several zeros $\rho_1,\ldots,\rho_m$ satisfy $\eta_k:=\beta_k-\sigma>0$,
the lower bound is dominated by the zero with maximal~$\eta_k$.
Indeed, each $\rho_k$ contributes independently to the derivative penalty
via Lemma~\ref{lem:derivative-penalty}.  Cross terms between distinct zeros
$\rho_k\ne\rho_\ell$ involve integrands of the form
$e^{(\eta_k+\eta_\ell)t}\cos(\gamma_k t)\cos(\gamma_\ell t)
=\tfrac12 e^{(\eta_k+\eta_\ell)t}
\bigl[\cos((\gamma_k-\gamma_\ell)t)+\cos((\gamma_k+\gamma_\ell)t)\bigr]$.
On a block $I_j$ of length $\Delta\asymp 1/\log T$, if $|\gamma_k-\gamma_\ell|\gg 1/\Delta$
the oscillatory integral contributes $O(1/|\gamma_k-\gamma_\ell|)$
by the Riemann--Lebesgue lemma;
by standard zero-spacing estimates~\cite{montgomery-pair}, consecutive zeros satisfy
$|\gamma_k-\gamma_\ell|\gg 1/\log T\asymp\Delta$,
so these cross terms are $O(1)$ per block and contribute at most $O(T\log T)$
after summation---subexponential relative to the dominant term.
Consequently,
\[
F_\lambda(H_\sigma;T,\Delta)
\;\gg\;
\max_{1\le k\le m}
\frac{\eta_k^{2}}{|\rho_k|^{2\alpha}}\,
\min\!\Bigl\{\Delta,\tfrac{1}{|\gamma_k|}\Bigr\}\,
\frac{e^{2\eta_k T}}{\eta_k\Delta},
\]
with the implied constant depending only on $(\alpha,\sigma)$ and the
canonical parameters.  For $T$ sufficiently large, the exponential with
$\eta^{\star}=\max_k\eta_k$ dominates all others.
Numerical confirmation via multi-zero injection appears in
Section~\ref{sec:multi-zero}.
\end{remark}

% ---------------------------------------------------------
\section{The Barrier Dichotomy and the Horocycle Conjecture}\label{sec:barrier}

\begin{proposition}[Barrier: proved direction]\label{prop:barrier}
If all zeros of $\zeta(s)$ satisfy $\beta\le\sigma$, then
$F_\lambda=O(T\log T\log\log T)$
\textup{(Theorem~\ref{thm:variance})}.
Conversely, if any zero satisfies $\beta>\sigma$, then
$F_\lambda\gg e^{2(\beta-\sigma)T}$
\textup{(Theorem~\ref{thm:bias})}.
\end{proposition}

\begin{proof}
The forward direction is Theorem~\ref{thm:variance} (proved in Part~1);
the contrapositive of the converse is Theorem~\ref{thm:bias} (proved above).
\end{proof}

The contrapositive of detection already yields the converse for \emph{finite}
configurations.

\begin{proposition}[Converse for finite off-line sets]\label{prop:finite-converse}
Suppose $\zeta$ has finitely many \textup{(}and at least one\textup{)} zeros with
$\beta>\sigma$, of maximal excess $\eta^\ast=\max(\beta-\sigma)>0$.  Then
$F_\lambda(H_\sigma;T,\Delta)\gg e^{2\eta^\ast T}$, so
$\sup_{T\ge T_0}F_\lambda/(T\log T\log\log T)=\infty$.  Hence if that supremum is
finite, $\zeta$ has no finite nonempty set of zeros with $\beta>\sigma$.
\end{proposition}

\begin{proof}
Apply Theorem~\ref{thm:bias} to a zero attaining $\eta^\ast$; the remaining
finitely many off-line zeros are absorbed by Remark~\ref{rmk:multiple-zeros}, and
the on/left zeros contribute the polynomial residual of
Step~\ref{step:residual-variance}.  Since $e^{2\eta^\ast T}$ outgrows every
polynomial, the normalized supremum diverges.
\end{proof}

\begin{conjecture}[Horocycle Conjecture]\label{conj:horocycle}
For every $\sigma\in[\tfrac12,1)$, if
$\sup_{T\ge T_0} F_\lambda(H_\sigma;T,\Delta)/(T\log T\log\log T)<\infty$,
then all zeros of $\zeta(s)$ satisfy $\beta\le\sigma$.
\end{conjecture}

\noindent
By Proposition~\ref{prop:finite-converse} the conjecture holds whenever the zeros
with $\beta>\sigma$ are finite in number.  It is \emph{not} merely the
contrapositive of Theorem~\ref{thm:bias}: the detection bound is proved only for a
single off-line zero against an on/left background, so its open content lies
precisely in the configurations that method cannot reach --- infinite off-line
families with $\beta_n\uparrow\beta^\ast$ \emph{unattained}, and off-line zeros
with $\beta-\sigma\to0$ (excluded from the threshold guarantee $\eta\ge c/\log T$
of Lemma~\ref{lem:threshold}).  Equivalently it asks for a uniform \emph{lower}
bound on $F_\lambda(H_\sigma)$ over arbitrary off-line configurations,
complementary to the upper bounds of Appendix~\ref{app:A}.

\subsection{A rigidity reformulation of the open case}\label{subsec:rigidity}

The obstruction is one of \emph{rigidity}: whether zeros accumulating at the
spectral edge can conspire, through cancellation, to suppress the growth rate of
$F_\lambda$.  Write $\beta^\ast=\sup\{\beta:\zeta(\beta+i\gamma)=0\}$.  We first
dispose of the \emph{attained} case, which already strengthens
Proposition~\ref{prop:finite-converse}.

\begin{proposition}[Attained edge]\label{prop:attained-edge}
If $\beta^\ast>\sigma$ and the supremum is attained \textup{(}by finitely many
zeros\textup{)}, then $F_\lambda(H_\sigma;T,\Delta)\gg e^{2(\beta^\ast-\sigma)T}$,
even when infinitely many further zeros lie in $(\sigma,\beta^\ast)$.
\end{proposition}

\begin{proof}[Sketch]
Write $F_\lambda=\|R\|_{\mathrm{pen}}^2$ with $R=(\mathrm{Id}-P)H_\sigma$ and
$\|\cdot\|_{\mathrm{pen}}^2=\sum_j(\|\cdot\|_{L^2(I_j)}^2+\lambda\|\partial_t\cdot\|_{L^2(I_j)}^2)$.
Let $w=(\mathrm{Id}-P)h^\ast$, where $h^\ast=\sum h_\rho$ runs over the finitely
many zeros with $\beta=\beta^\ast$; these have distinct, well-separated ordinates,
so $\|w\|_{\mathrm{pen}}^2\asymp e^{2(\beta^\ast-\sigma)T}$ (the single-zero
estimate of Step~\ref{step:derivative-lb}).  By Cauchy--Schwarz,
$F_\lambda\ge\langle R,w\rangle_{\mathrm{pen}}^2/\|w\|_{\mathrm{pen}}^2$.  The
pairing $\langle R,w\rangle_{\mathrm{pen}}=\sum_\rho\langle(\mathrm{Id}-P)h_\rho,w\rangle_{\mathrm{pen}}$
equals $\|w\|_{\mathrm{pen}}^2$ plus cross terms with the lower zeros, each of
exponential rate $\eta_\rho+\eta^\ast<2\eta^\ast$ and hence
$o(e^{2(\beta^\ast-\sigma)T})$.  Thus
$\langle R,w\rangle_{\mathrm{pen}}\asymp e^{2(\beta^\ast-\sigma)T}$ and
$F_\lambda\gg e^{4(\beta^\ast-\sigma)T}/e^{2(\beta^\ast-\sigma)T}=e^{2(\beta^\ast-\sigma)T}$.
\end{proof}

The single unresolved case is $\beta^\ast$ \emph{unattained}: zeros
$\beta_n\uparrow\beta^\ast$ with none achieving the supremum.  Then no finite top
test field $w$ exists, and the cross terms from zeros with $\eta_\rho$ arbitrarily
close to $\eta^\ast$ share the diagonal's exponential rate, so cancellation is not
excluded by the elementary argument.

\begin{remark}[Rigidity formulation]\label{rmk:rigidity}
Since $F_\lambda(H_\sigma;\,\cdot\,)\ge0$ as a function of $T$, the
Laplace-transform singularity theorem for nonnegative integrands
(Pringsheim's theorem; \cite[Ch.~II]{widder}) identifies the abscissa of
convergence of $\int_0^\infty e^{-uT}F_\lambda\,dT$ with a real singularity,
equal to the exponential growth rate of $F_\lambda$.  The \emph{diagonal} part of
$F_\lambda$ has its rightmost singularity at $u=2(\beta^\ast-\sigma)$ with
nonnegative residues; the open case is exactly the assertion that the
off-diagonal part cannot cancel this singularity and move the abscissa left.
Precisely, the reduction is $u_c=2\,\limsup_{t\to\infty}\tfrac1t\log|H_\sigma'(t)|$,
so the question is whether the growth (boundedness) abscissa of the exponential
sum $H_\sigma'$ equals its absolute-convergence abscissa $\eta^\ast=\beta^\ast-\sigma$.
For general exponential sums these can differ (a Bohr phenomenon): the
term-by-term pole locations of the formal $M(u)$ do not control the continued
singularity, and near-edge cancellation could in principle lower $u_c$ below
$2\eta^\ast$ --- but, by Proposition~\ref{prop:attained-edge}, \emph{only when
$\beta^\ast$ is unattained}.  Zero-density bounds give only the trivial upper
bound $u_c\le2\eta^\ast$; excluding the downward cancellation is the genuine
rigidity content.  This
is a \emph{rigidity} statement --- the growth rate of $F_\lambda$ is conjecturally
the spectral edge $2(\beta^\ast-\sigma)_+$, undisturbed by cancellation among
edge-accumulating zeros --- and is the analytic counterpart of the
curvature/horocycle rigidity of Part~4, morally parallel to the rigidity of
unipotent and horocycle flows, where no exotic invariant structures intervene.
We isolate it as the precise content of the open case: \emph{bounded energy
rigidly confines the spectrum to $\beta\le\sigma$.}
\end{remark}

\begin{remark}[A Gram-matrix (Ingham) formulation]\label{rmk:gram}
The rigidity is a quantitative linear-independence statement, controlled by the
Gram matrix of the edge exponentials.  Writing
$e_\rho(t)=e^{(\eta_\rho+i\gamma_\rho)t}$ (with $\eta_\rho$ near $\eta^\ast$),
cancellation of near-edge contributions is exactly degeneracy of their Gram
matrix $G(T)_{\rho\rho'}=\int_0^T e_\rho(t)\,\overline{e_{\rho'}(t)}\,dt$, and a
spectral gap precludes it.  (This is the same Gram-nondegeneracy that governs
correctability in the Knill--Laflamme conditions for quantum
codes~\cite{knill-laflamme}, where diagonalizing the Hermitian
$c_{ab}=PE_a^\dagger E_bP$ splits the errors, by polar decomposition, into a
correctable and a vanishing part.)  A direct
computation shows the \emph{normalized} Gram matrix is essentially independent
of $T$:
\begin{equation}
\widetilde G_{\rho\rho'}
=\frac{\langle e_\rho,e_{\rho'}\rangle}{\|e_\rho\|\,\|e_{\rho'}\|}
\;\xrightarrow{\,T\to\infty\,}\;
\frac{2\sqrt{\eta_\rho\,\eta_{\rho'}}}{\sqrt{(\eta_\rho+\eta_{\rho'})^2+(\gamma_\rho-\gamma_{\rho'})^2}}
\;\le\;\frac{2\eta^\ast}{|\gamma_\rho-\gamma_{\rho'}|}
\qquad(\rho\ne\rho'),
\label{eq:gram-cos}
\end{equation}
with unit diagonal; as a genuine cosine it never exceeds $1$, and the right-hand
bound is the well-separated form $|\gamma_\rho-\gamma_{\rho'}|\gg\eta^\ast$.  Thus
the \emph{angle} between two edge contributions is fixed by their ordinate gap:
cancellation (near-degeneracy of $\widetilde G$) requires edge zeros with
\emph{close ordinates}, whereas ordinate separation forces near-orthogonality and
hence rigidity.  Concretely, since the off-diagonal entries are bounded by
$2\eta^\ast/|\gamma_\rho-\gamma_{\rho'}|$, the Schur test makes $\widetilde G$
positive definite with a spectral gap as soon as
$\sup_\rho\sum_{\rho'\ne\rho}2\eta^\ast/|\gamma_\rho-\gamma_{\rho'}|<1$, in which
case $F_\lambda\gg e^{2(\eta^\ast-\varepsilon)T}$ and no cancellation occurs (a
weighted, finite-$T$ analogue of Ingham's inequality~\cite{ingham-ineq}).  Thus
the unattained-edge rigidity is, like Hypothesis~(S) of Appendix~\ref{app:A}, an
instance of \emph{non-clustering of ordinates} --- though at much larger heights
and in a far stronger (uniform, over an unbounded range) form than the trivial
low-height separation that suffices for the variance bound.  We do not establish
the Schur bound
unconditionally---over an unbounded range of heights the row sums need not be
finite---so the rigidity is \emph{reduced to}, but not closed by, ordinate
separation.
\end{remark}

% ---------------------------------------------------------
\section{Heuristic and Numerical Consistency}\label{sec:numerics}

In the threshold regime $\eta=1/\log T$, $e^{2\eta T}=e^{2T/\log T}$
eventually dominates any polynomial, sharply separating variance and bias behaviour.
The numerical experiments in Part~3 (Odlyzko zeros plus synthetic off-line injections)
show:
(i) polynomial growth under $\beta\le\sigma$, consistent with
Theorem~\ref{thm:variance}; and
(ii) exponential growth with empirical slope matching $2\eta$ to within $3\times10^{-4}$ relative error.

\paragraph{Finite-$T$ caveat.}
For very small $\eta$, the asymptotic slope appears only once $T\gg1/\eta$;
the observed trend is monotone and converges to the theoretical limit.

% ---------------------------------------------------------
\section{Geometric Preview}\label{sec:geom-preview}

The exponential growth proven above is the analytic counterpart of a
radial escape in a variable-curvature model space:
trajectories with $\beta>\sigma$ correspond to $u=e^{\eta t}>1$
and exhibit geodesic-like expansion.
Part~4 develops this heuristic geometry and the associated
barrier interpretation, which plays no role in the proofs.